\chapter{Introducción y Objetivos}
\label{chap:introduccion}

Este trabajo práctico tiene como propósito principal el empleo de métodos de medición e instrumentos
sencillos para la determinación de las características de una red de dos puertos, centrándose en el análisis
de amplificadores de baja frecuencia. Se estudian métodos prácticos para medir impedancias y ganancias,
además del impacto del uso de escalas logarítmicas y los efectos del lazo de realimentación negativa.

\section{Objetivos}
\begin{itemize}
    \item Determinar los parámetros fundamentales de un amplificador: impedancia de entrada ($Z_i$),
          impedancia de salida ($Z_o$), ganancia de tensión ($A_v$) y ganancia de potencia ($A_p$).
    \item Adquirir destreza en el uso de instrumentos de medición con escalas graduadas en decibeles ($dB$),
          comprendiendo el significado físico de las unidades $dBu$ y $dBm$.
    \item Emplear el osciloscopio como herramienta principal para efectuar el análisis dinámico y la
          medición de parámetros temporales y de amplitud en sistemas de amplificación.
    \item Evaluar experimentalmente el efecto de la realimentación negativa sobre la ganancia, la respuesta
          en frecuencia (ancho de banda), la distorsión y las impedancias de entrada y salida.
    \item Expresar los resultados obtenidos en las mediciones acompañados de su correspondiente grado de
          incertidumbre experimental.
\end{itemize}

\section{Introducción Teórica}

La caracterización de un amplificador lineal operando en baja frecuencia se puede realizar mediante un modelo
de circuito equivalente de dos puertos. Este modelo simplificado queda unívocamente determinado por la
impedancia de entrada ($Z_i$), la impedancia de salida ($Z_o$) y la ganancia de tensión en vacío ($A_{v0}$).

\subsection{Medición de Impedancias por el Método de Reducción a la Mitad}
La impedancia de salida se determina indirectamente mediante un método de sustitución. Midiendo la tensión de
salida en vacío ($V_s$) y luego conectando una carga variable hasta que dicha tensión disminuya a la mitad
($V_s' = V_s/2$), se deduce que el valor de la carga adaptada es idéntico al de la impedancia interna de salida.
Para la impedancia de entrada, se intercala un resistor variable en serie con el generador y se repite el proceso
de reducción a la mitad de la tensión de salida.

\subsection{Uso de la Escala en Decibeles}
El decibel ($dB$) es una unidad logarítmica que expresa la relación entre dos magnitudes. En audio y telecomunicaciones,
el voltímetro de CA está calibrado de manera que $0\dBu$ equivale a una tensión de referencia de $0,775\V$. Cuando esta
tensión se aplica sobre una carga patrón de $600\ohm$, la potencia disipada es de $1\mW$, lo cual define el $0\dBm$.
Si la impedancia de carga difiere de la de referencia ($600\ohm$), debe aplicarse un término corrector:
\begin{equation}
    dBm = dBu + 10 \cdot \log_{10} \left( \frac{600\ohm}{R_x} \right)
\end{equation}
donde $R_x$ es la resistencia sobre la cual se realiza la medición.

\subsection{Efecto de la Realimentación Negativa}
La realimentación negativa consiste en tomar una muestra de la señal de salida y restarla en la entrada.
En la configuración de muestreo de tensión y mezcla en serie, la ganancia a lazo cerrado ($G$) disminuye respecto
de la ganancia a lazo abierto ($A$) según:
\begin{equation}
    G = \frac{A}{1 + A \cdot B}
\end{equation}
donde $B$ es la ganancia del lazo de realimentación. La realimentación estabiliza la ganancia, reduce la
distorsión armónica, aumenta la impedancia de entrada y reduce la impedancia de salida a lazo cerrado ($Z_o'$):
\begin{equation}
    Z_o' = \frac{Z_o}{1 + A \cdot B} = \frac{Z_o \cdot G}{A}
\end{equation}
Además, incrementa el ancho de banda del amplificador a expensas de reducir la ganancia en frecuencias medias.
El ancho de banda se puede evaluar mediante el barrido senoidal en frecuencia o de forma indirecta midiendo el
tiempo de crecimiento ($t_c$) ante un escalón de tensión (onda cuadrada):
\begin{equation}
    AB \approx \frac{0,35}{t_c}
\end{equation}
Este método permite una caracterización rápida del comportamiento dinámico del circuito.
